\documentclass[12pt, a4paper]{article}

\usepackage[utf8]{inputenc}
\usepackage[T2A]{fontenc}
\usepackage[russian]{babel}

\usepackage[left=2.5cm, right=1.5cm, top=2cm, bottom=2cm]{geometry}

\usepackage{indentfirst}
\usepackage{amsmath}
\usepackage{graphicx}
\usepackage{hyperref}
\usepackage{listings}
\usepackage{xcolor}
\usepackage{booktabs}
\usepackage{float}

\definecolor{codegreen}{rgb}{0,0.6,0}
\definecolor{codegray}{rgb}{0.5,0.5,0.5}
\definecolor{codepurple}{rgb}{0.58,0,0.82}
\definecolor{backcolour}{rgb}{0.95,0.95,0.92}

\lstdefinestyle{mystyle}{
    backgroundcolor=\color{backcolour},
    commentstyle=\color{codegreen},
    keywordstyle=\color{magenta},
    numberstyle=\tiny\color{codegray},
    stringstyle=\color{codepurple},
    basicstyle=\ttfamily\footnotesize,
    breakatwhitespace=false,
    breaklines=true,
    captionpos=b,
    keepspaces=true,
    numbers=left,
    numbersep=5pt,
    showspaces=false,
    showstringspaces=false,
    showtabs=false,
    tabsize=4
}

\lstset{style=mystyle}

\begin{document}
\begin{center}
\large
МОСКОВСКИЙ АВИАЦИОННЫЙ ИНСТИТУТ\\ (НАЦИОНАЛЬНЫЙ ИССЛЕДОВАТЕЛЬСКИЙ УНИВЕРСИТЕТ)

\vspace{20pt}

Институт №8 <<Компьютерные науки и прикладная математика>>

Кафедра 806 <<Вычислительная математика и программирование>>
\end{center}

\vspace{60pt}

\begin{center}
\bfseries
\large
Лабораторная работа № 1 по курсу дискрeтного анализа: сортировка за линейное время

\vspace{54pt}

\end{center}

\vfill

\begin{flushright}
\large
\begin{tabular}{rl}
Выполнил: & Иванов Иван \\
Группа: & М8О-207БВ-24 \\
\end{tabular}
\end{flushright}

\vspace{92pt}

\begin{center}
\large
Москва, \the\year
\end{center}
\newpage

\section{Условие}

\subsection{Общая постановка задачи}
Требуется разработать программу, осуществляющую ввод пар «ключ-значение», их упорядочивание по возрастанию ключа указанным алгоритмом сортировки за линейное время и вывод отсортированной последовательности.

\subsection{Вариант задания: 4-1}
Поразрядная сортировка.\\
Тип ключа: даты в формате DD.MM.YYYY, например 1.1.1, 1.9.2009, 01.09.2009, 31.12.2009.\\
Тип значения: строки фиксированной длины 64 символа, во входных данных могут встретиться строки меньшей длины, при этом строка дополняется до 64-х нулевыми символами, которые не выводятся на экран.

\section{Метод решения}

Для решения задачи была реализована поразрядная сортировка (Radix Sort), которая обеспечивает линейную временную сложность O($n \times k$), где n - количество элементов, k - количество разрядов в ключе.

\subsection{Преобразование ключа}
Даты в формате DD.MM.YYYY преобразуются в целочисленный ключ формата YYYYMMDD для возможности сравнения:
$Key = Year \times 10000 + Month \times 100 + Day$

\subsection{Алгоритм сортировки}
Поразрядная сортировка работает следующим образом:\\
Находим максимальный ключ для определения количества разрядов. Для каждого разряда (от младшего к старшему) подсчитываем количество элементов с каждой цифрой в текущем разряде, вычисляем позиции элементов в результирующем массиве, распределяем элементы по позициям, копируем отсортированные данные обратно.

\subsection{Архитектура программы}
Программа состоит из:
\begin{itemize}
    \item Структуры \texttt{TRecord} для хранения пары ключ-значение
    \item Структуры \texttt{TVector} для динамического массива записей
    \item Функции управления вектором (\texttt{InitTVector}, \texttt{PushTVector}, \texttt{FreeTVector}, \texttt{RecapacityTVector})
    \item Функция парсинга даты \texttt{ParseDate}
    \item Функция поразрядной сортировки \texttt{RadixSort}
\end{itemize}

\section{Описание программы}

\subsection{Структура файлов}
\texttt{main.cpp} — файл программы

\subsection{Основные типы данных}
\begin{lstlisting}[language=C++, caption={Структуры}]
struct TRecord {
    int Key;
    char Date[DATE_SIZE];
    char Value[VALUE_SIZE];
};

struct TVector {
    TRecord *Data;
    int Size;
    int Capacity;
};
\end{lstlisting}

\subsection{Основные функции}
\begin{itemize}
    \item \texttt{InitTVector} - Инициализация вектора
    \item \texttt{PushTVector} - Добавление элемента в вектор
    \item \texttt{RecapacityTVector} - Изменение емкости вектора
    \item \texttt{FreeTVector} - Уничтожение вектора
    \item \texttt{ParseDate} - Парсинг даты в ключ
    \item \texttt{GetMaxKey} - Поиск максимального ключа
    \item \texttt{RadixSort} - Поразрядная сортировка
\end{itemize}

\section{Дневник отладки}
Возникли проблемы с парсингом дат с однозначными числами. Было решено использовать sscanf с форматом\%d вместо \%02d. Была упощена вероятность того, что пользователь введет дату в неверном формате, была добавлена проверка на DATE\_FIELDS\_COUNT.

\section{Недочеты}

Программа не обрабатывает некорректные форматы дат (например, 32.13.2024). Отсутствует проверка на переполнение при выделении памяти для очень больших массивов.

\section{Выводы}

В ходе выполнения лабораторной работы была реализована программа сортировки пар «ключ-значение» с использованием поразрядной сортировки. Реализованный алгоритм может применяться для cортировки больших объемов данных с целочисленными ключами, сортировки дат, идентификаторов, телефонных номеров. Также алгоритм подходит для решения задач, где требуется устойчивая сортировка за линейное время. Основная сложность заключалась в работе с динамической памятью без использования STL, задача свелась к применению парадигм языка C. Также проблемы возникли при обработке различных форматов записи дат.
\end{document}
